\centerline{\bf \Huge Комбинаторика. База.}


\textbf{Правило произведения}

Пусть объект А можно выбрать $n$ способами и после каждого такого выбора объект В можно выбрать $m$ способами. Тогда выбор пары ($\mathrm{A}, \mathrm{B}$) можно осуществить $n m$ способами.

\textit{Пример 1.}
Пусть $p_1, p_2, \ldots, p_n$ - различные простые числа; $k_1, k_2, \ldots, k_n$ - целые неотрицательные числа. Сколько делителей у числа $a=p_1^{k_1} p_2^{k_2} \ldots p_n^{k_n}$ ?

\textbf{Правило суммы}

Пусть некоторый объект А можно выбрать $n$ различными способами, а другой объект В можно выбрать $\boldsymbol{m}$ способами. Тогда существует $\boldsymbol{n}+\boldsymbol{m}$ способов выбрать либо объект А, либо объект В.

\textbf{Правило дополнения}

Если $A = B \cup C$, и множества $B$ и $C$ не пересекаются, то мощность множества $A$ вычисляется по формуле $|A| = |B| + |C|$ (в общем случае формула выглядит так $|A| = |B| + |C| - |B\cap C|$). В данном случае множество $B$ дополняет множество $C$ до множества $A$. Выглядит очевидно, но зачастую в задачах требуют посчитать мощность какого-то неудобного множества, у которого удобное дополнение.


\textit{Пример 2.} Сколько трёхзначных чисел содержат в своей записи хотя бы одну чётную цифру?


Способ поставить в ряд $n$ различных объектов называется \textbf{перестановкой} этих объектов.

\textit{Пример 3.} Сколько слов (не обязательно осмысленных) можно получить, переставляя буквы в слове ДУБ?

Число перестановок из $n$ элементов обозначается $P_n$. Вычисление $P_n$ аналогично решению примера: на первое место можно поставить любой из $n$ объектов, на второе - любой из $n-1$ оставшихся, $\ldots$, на предпоследнее - любой из двух оставшихся и, наконец, последнее место останется для последнего объекта. Значит, $P_n=n \cdot(n-1) \cdot \ldots \cdot 2 \cdot 1$.

Произведение всех натуральных чисел от 1 до $n$ называется факториалом числа $n$ и обозначается $n!$ (читается «эн факториал»): $n!=1 \cdot 2 \cdot 3 \cdot \ldots \cdot n$. Для удобства считают $0!=1!=1$. Итак,

$$
P_n=n!.
$$

\begin{flushright}
        \vspace{-10mm}
        (1)
\end{flushright}

\textit{Пример 4.} Сколькими способами можно переставить буквы в слове ЗАДАЧА?

\textit{Пример 5.} Сколькими способами можно переставить буквы в слове КОЛОКОЛ?

Эти два примера иллюстрируют понятие перестановки с повторениями. Рассмотрим общий случай. Пусть есть $k_1$ предметов 1 -го типа, $k_2-2$-го, $\ldots, k_m-m$-го типа. Способ поставить эти $n=k_1+k_2+\ldots+k_m$ предметов в ряд называется перестановкой из $n$ элементов с заданным числом повторений $k_1, k_2, \ldots, k_m$. Число таких перестановок обозначается $\bar{P}_{k_1, k_2, \ldots, k_m}$ и вычисляется по формуле

$$
\bar{P}_{k_1, k_2, \ldots, k_m}=\frac{\left(k_1+\ldots+k_m\right)!}{k_{1}!k_{2}!\ldots k_{m}!},
$$
\begin{flushright}
        \vspace{-10mm}
        (2)
\end{flushright}

вывод которой аналогичен решению только что разобранных примеров.

\textit{Пример 6.} На школьном собрании выбирают старосту, помощника старосты и дежурного (причём это должны быть разные люди). Сколькими способами это можно сделать, если в классе учится 30 человек?

Пусть теперь даны $n$ объектов и число $k$ от 0 до $n$. Размещением (без повторений) из $n$ элементов по $k$ называется упорядоченная выборка длины $k$ из этих $n$ объектов, где каждый объект встречается по одному разу. Число таких размещений обозначается $A_n^k$ и, по аналогии с только что разобранным примером, вычисляется так:

$$
A_n^k=n(n-1) \ldots(n-(k-1))=\frac{n!}{(n-k)!} .
$$
\begin{flushright}
        \vspace{-10mm}
        (3)
\end{flushright}

В частности, при $k=n$ получаем число перестановок: $A_n^n=P_n$; при $k=0$ получаем единственную пустую выборку: $A_n^0=1$.

Сочетания отличаются от размещений тем, что в них не важен порядок выбранных объектов. Иными словами, сочетанием из $n$ по $k$ называется неупорядоченная выборка, содержащая какие-то $k$ из этих $n$ объектов, или $k$ элементное подмножество данного $n$-элементного множества объектов. Число сочетаний из $n$ по $k$ обозначается $C_n^k$ или $\binom{n}{k}$ (мы будем придерживаться первого обозначения).

Поскольку неупорядоченную $k$-элементную выборку можно упорядочить $P_k=k!$ способами, число $A_n^k$ размещений в $k!$ раз больше числа $C_n^k$ сочетаний, откуда с учётом (3) получаем формулу для числа сочетаний:

$$
C_n^k=\frac{n!}{(n-k)!k!} .
$$
\begin{flushright}
        \vspace{-10mm}
        (4)
\end{flushright}


Например, трёх человек из тридцати можно выбрать $\frac{30 \cdot 29 \cdot 28}{3 \cdot 2 \cdot 1}$ способами.

Сравнивая полученную формулу для числа сочетаний с формулой (2), можно заметить, что $C_n^k=\bar{P}_{k, n-k}$. Это не случайно. Действительно, каждому сочетанию из $n$ по $k$ можно сопоставить строку длины $n$ из нулей и единиц, в которой на $i$-м месте стоит единица, если $i$-й элемент входит в сочетание, и ноль, если не входит. Значит, $C_n^k$ равно числу строк из $k$ единиц и $n-k$ нулей, а это и есть $\bar{P}_{k, n-k}$. Это рассуждение является ещё одним способом

\textbf{Бином Ньютона.} Раскроем скобки в выражении

$$
(a+b)^n=\underbrace{(a+b)(a+b) \ldots(a+b)}_{n \text { скобок }} .
$$


Поскольку из каждой скобки можно взять любую из двух букв, а всего скобок $n$, при их раскрытии получится сумма $2^n$ слагаемых, каждое из которых имеет вид $a^{n-k} b^k$ для некоторого $k$ от 0 до $n$. Приведём подобные. Коэффициент при слагаемом $a^{n-k} b^k$ (при фиксированном $k$ ) равен числу способов взять из $n$ скобок $k$ букв $b$, т. е. $C_n^k$. Итак, получаем формулу

$$
(a+b)^n=a^n+C_n^1 a^{n-1} b+\ldots+C_n^k a^{n-k} b^k+\ldots+C_n^{n-1} a b^{n-1}+b^n
$$
\begin{flushright}
        \vspace{-10mm}
        (5)
\end{flushright}

называемую формулой \textit{бинома Ньютона}. Поэтому числа $C_n^k$ называются ещё и \textit{биномиальными коэффициентами}.

\textbf{Свойства биномиальных коэффициентов.} Следующие свойства (а также тождества с биномиальными коэффициентами, оставленные в качестве задач в конце параграфа) могут быть доказаны при помощи формулы (4). Однако многие из них можно доказать гораздо изящнее при помощи комбинаторных рассуждений или бинома Ньютона.

\textit{Симметрия}: $C_n^k=C_n^{n-k}$. Действительно, всякая $k$-элементная выборка из $n$ элементов однозначно определяет выборку из $n-k$ остальных элементов. Так же просто это видно и из формулы (4), в которой при замене $k$ на $n-k$ меняются местами факториалы в знаменателе.

\textit{Основное биномиальное тождество}: $C_{n+1}^k=C_n^{k-1}+C_n^k$. Возьмём $n$ белых и один чёрный шар и посчитаем, сколькими способами можно выбрать из них $k$ шаров. С одной стороны, это число равно $C_{n+1}^k \cdot \mathrm{C}$ другой стороны, все выборки можно разделить на две группы: содержащие чёрный шар (таких $C_n^{k-1}$ ) и не содержащие (таких $C_n^k$ ).

\textit{Сумма биномиальных коэффициентов}: $C_n^0+C_n^1+\ldots+C_n^n=2^n$. Эту формулу можно получить подстановкой $a=b=1$ в бином Ньютона. Есть и другой - комбинаторный - способ доказательства. Посчитаем число подмножеств $n$-элементного множества. С одной стороны, все подмножества можно классифицировать по количеству элементов; т. к. $k$-элементных подмножеств ровно $C_n^k$, всего подмножеств будет $C_n^0+C_n^1+\ldots+C_n^n$. С другой стороны, любое подмножество определяется тем, какие элементы в него входят, а какие нет. Поскольку для каждого из $n$ элементов исходного множества есть две возможности: входить или не входить в подмножество, подмножеств всего $2^n$.